\chapter{Conclusion générale}

Ce projet de fin d'études a permis de concevoir et de réaliser \textbf{MoroccoCheck}, une plateforme touristique communautaire articulée autour d'une \textbf{API REST}, d'une application \textbf{Flutter} et d'une interface d'\textbf{administration web}. Les objectifs initiaux — catalogue de sites, \textbf{vérification par géolocalisation}, \textbf{avis}, \textbf{gamification}, rôles métiers et \textbf{modération} — sont couverts par une base logicielle cohérente et documentée.

Sur le plan technique, l'équipe a consolidé des compétences en développement \textbf{backend} (Node.js, MySQL, sécurité), \textbf{mobile} (Flutter, intégration API) et \textbf{frontend web} (React), ainsi qu'en modélisation (\textbf{UML}), gestion de version (\textbf{Git}) et intégration continue (\textbf{GitHub Actions}).

Les perspectives à moyen terme portent sur la généralisation du déploiement en environnement de production, l'élargissement des tests automatisés et l'enrichissement progressif des fonctionnalités (recommandations, analytics, expérience multilingue, etc.), dans la continuité des axes identifiés dans l'analyse du projet.

Ce travail démontre qu'une approche \textbf{itérative}, fondée sur une architecture claire et des besoins bien identifiés, permet de livrer un système complexe tout en gardant une trajectoire d'amélioration maîtrisée.
